% !TITLE 永遇乐·京口北固亭怀古
% !AUTHOR 辛弃疾
% !DYNASTY 两宋
% !GENRE 词
% !ORDER 15

\begin{poem}
千古江山，英雄无觅，孙仲谋\textsuperscript{1}处。
舞榭歌台\textsuperscript{2}，风流总被，雨打风吹去。
斜阳草树，寻常巷陌，人道寄奴\textsuperscript{3}曾住。
想当年，金戈铁马\textsuperscript{4}，气吞万里如虎。

元嘉草草\textsuperscript{5}，封狼居胥\textsuperscript{6}，赢得仓皇北顾\textsuperscript{7}。
四十三年\textsuperscript{8}，望中犹记，烽火扬州路。
可堪回首，佛狸祠下\textsuperscript{9}，一片神鸦社鼓\textsuperscript{10}。
凭谁问，廉颇老矣，尚能饭否\textsuperscript{11}？
\end{poem}

\begin{annotations}
\notetext{1}{孙仲谋：孙权（182-252），字仲谋，三国吴国开国皇帝，曾在京口（今镇江）建都。辛弃疾登京口北固亭，首先想到的就是这位英雄。}
\notetext{2}{舞榭歌台：歌舞楼台，代指昔日的繁华风流。再辉煌的事业也被时光的雨打风吹去。}
\notetext{3}{寄奴：南朝宋武帝刘裕的小名。刘裕出身贫寒，曾住在京口寻常的巷子里，后来成为一代雄主。}
\notetext{4}{金戈铁马：金属长戈和披铁甲的战马。刘裕统率精锐之师北伐，气势如吞噬万里的猛虎。}
\notetext{5}{元嘉草草：元嘉是宋文帝的年号。宋文帝草率出兵北伐，大败而归。辛弃疾以此告诫南宋朝廷：北伐不能草率行事。}
\notetext{6}{封狼居胥：在狼居胥山（今蒙古境内）筑坛祭天，庆贺胜利。汉武帝时霍去病大败匈奴，封狼居胥山——这是古代武将的最高荣誉。}
\notetext{7}{仓皇北顾：仓皇地回头看北方追兵。指元嘉北伐失败，反被北魏追击。}
\notetext{8}{四十三年：辛弃疾从率众归宋到写此词，正好四十三年。他犹记得当年在扬州一带与金兵作战的烽火。}
\notetext{9}{佛狸祠：佛狸是北魏太武帝拓跋焘的小名，他曾追击宋军到长江北岸的瓜步山，在那里建行宫，后被改作祠庙。佛（bì），通“狴”。}
\notetext{10}{神鸦社鼓：神鸦吃祭品的乌鸦，社鼓是社日祭神的鼓声。当年的战场成了祠庙，老百姓在那里祭神击鼓，好像什么也没发生过——这对辛弃疾来说是难以忍受的遗忘。}
\notetext{11}{廉颇老矣，尚能饭否：廉颇是战国赵国名将，晚年被赵王怀疑，使者去看他，廉颇在使者面前吃了一斗米、十斤肉，表示还能打仗，但使者受贿说廉颇老不中用了。辛弃疾以廉颇自比——我也老了，还有人愿意用我吗？}
\end{annotations}

\begin{appreciation}
《永遇乐·京口北固亭怀古》写于辛弃疾六十五岁那年，朝廷要北伐，起用他做镇江知府。他登上北固亭，写下两宋怀古词里最沉痛的一首之一。

京口就是今天的镇江，曾经是孙权建都的地方，所以开篇就是"千古江山，英雄无觅，孙仲谋处"——江山还在，孙权那样的英雄却找不到了。"舞榭歌台，风流总被雨打风吹去"——再辉煌的歌舞楼台也经不住时光。这是怀古词里最老的话，辛弃疾却马上落到具体处："斜阳草树，寻常巷陌，人道寄奴曾住"——寄奴是刘裕的小名，南朝宋的开国皇帝，出身贫寒，住京口一条普通巷子。"想当年，金戈铁马，气吞万里如虎"——一个"虎"字，气势呼之欲出。

上阕全是英雄赞歌，下阕笔锋一转，全是教训。"元嘉草草，封狼居胥，赢得仓皇北顾"——宋文帝刘义隆想学霍去病封狼居胥山，草率北伐，结果大败，仓皇南逃。这半阕是说给朝廷听的：北伐我赞成，但别学元嘉草草。后来的开禧北伐果然大败——他说中了。

"四十三年，望中犹记，烽火扬州路"——他南归整整四十三年，金兵打到扬州的那场战火，他一直记着。可是如今呢？"佛狸祠下，一片神鸦社鼓"——佛狸是当年追击宋军到长江边的北魏皇帝的小名，如今祠庙香火鼎盛，百姓祭神敲鼓，大概没人记得供的是谁。对辛弃疾来说，最可怕的不是失败，是遗忘：战场变成了庙会，仇恨变成了年例。

最后一句是全词的落点："凭谁问，廉颇老矣，尚能饭否"——廉颇老了，还能吃一斗米十斤肉、披甲上马，可探病的人收了钱，回去说他"坐一会儿就拉了三趟"。辛弃疾问：还有人来问我一句吗？没有人。

我们读过苏轼的《念奴娇·赤壁怀古》：苏轼怀周瑜，怀到最后是"一樽还酹江月"，敬明月一杯，把自己放过了；辛弃疾怀孙权怀刘裕，怀到最后是"凭谁问"——放不下。曹操说"老骥伏枥，志在千里"，好歹是自己喊的；廉颇连被问一声都等不到。辛弃疾一辈子想做的，就是《无衣》里那个"与子同袍"的战士；可是同袍的人在北方，他在南方，那些年没有人和他一起"修我戈矛"。写完这首词不到两年，辛弃疾去世，至死没能再上马。

哥哥选这首词，是想让你记住"记得"这两个字。人一辈子会忘记很多事，没办法；但有些东西要攥住：家里老人的故事，你小时候住过的巷子，还有你自己从哪里来。别让它们像佛狸祠那样，变成一场热闹的、没人记得缘由的庙会。
\end{appreciation}
